%=========== ΛΥΣΗ - 64ο ΘΕΜΑ Β ===========
\begin{thema}{Β}
Δίνεται η συνάρτηση $f(x)=\ln(x^2+1)+ax^2+\beta,\ x\in\mathbb{R}$.

\begin{erwthma}
\item Επειδή η $C_f$ διέρχεται από την αρχή των αξόνων, έχουμε
\[
f(0)=0\Leftrightarrow \ln 1+\beta=0\Leftrightarrow \beta=0.
\]
Επίσης, αφού το $M(1,\ln2)$ είναι σημείο της $C_f$, ισχύει
\[
f(1)=\ln2\Leftrightarrow \ln2+a+\beta=\ln2\Leftrightarrow a=0.
\]
Άρα
\[
a=0,\quad \beta=0
\]
και επομένως
\[
f(x)=\ln(x^2+1).
\]

\item Θεωρούμε τη συνάρτηση
\[
g(x)=\ln(x^2+1)-x^2,\quad x\in\mathbb{R}.
\]
Τότε
\[
g'(x)=\frac{2x}{x^2+1}-2x
=2x\left(\frac{1}{x^2+1}-1\right)
=-\frac{2x^3}{x^2+1}.
\]
Επειδή $x^2+1>0$, το πρόσημο της $g'$ είναι το αντίθετο του προσήμου του $x$:
\begin{itemize}
\item $g'(x)>0\Leftrightarrow x<0$,
\item $g'(x)=0\Leftrightarrow x=0$,
\item $g'(x)<0\Leftrightarrow x>0$.
\end{itemize}
Άρα η $g$ είναι γνησίως αύξουσα στο $(-\infty,0]$ και γνησίως φθίνουσα στο $[0,+\infty)$. Συνεπώς παρουσιάζει ολικό μέγιστο στο $x=0$, με
\[
g(0)=\ln1-0=0.
\]
Επομένως $g(x)\leq0$ για κάθε $x\in\mathbb{R}$, δηλαδή
\[
\ln(x^2+1)\leq x^2,\quad x\in\mathbb{R}.
\]
Η ισότητα ισχύει μόνο για $x=0$.

\item Με $f(x)=\ln(x^2+1)$ έχουμε
\[
\lim_{x\to0^+}\frac{f(x)}{x}
=\lim_{x\to0^+}\frac{\ln(x^2+1)}{x}
\xlongequal[\text{\eng{DLH}}]{\frac{0}{0}}
\lim_{x\to0^+}\frac{2x}{x^2+1}=0.
\]

\item Για κάθε $x\in\mathbb{R}$ ισχύει $x^2+1\geq1$, άρα
\[
f(x)=\ln(x^2+1)\geq\ln1=0.
\]
Επειδή $f(0)=0$, η $f$ παρουσιάζει \textbf{ολικό ελάχιστο} στο $x=0$, με ελάχιστη τιμή $0$.

Για την κυρτότητα έχουμε
\[
f'(x)=\frac{2x}{x^2+1}
\quad\text{και}\quad
f''(x)=\frac{2(x^2+1)-4x^2}{(x^2+1)^2}
=\frac{2(1-x^2)}{(x^2+1)^2}.
\]
Επειδή $(x^2+1)^2>0$, το πρόσημο της $f''$ καθορίζεται από το $1-x^2$:
\begin{itemize}
\item $f''(x)>0\Leftrightarrow -1<x<1$,
\item $f''(x)=0\Leftrightarrow x=-1$ ή $x=1$,
\item $f''(x)<0\Leftrightarrow x<-1$ ή $x>1$.
\end{itemize}
Άρα η $f$ είναι \textbf{κοίλη} στα $(-\infty,-1]$, $[1,+\infty)$ και \textbf{κυρτή} στο $[-1,1]$. Επειδή η κυρτότητα αλλάζει στα $x=-1$ και $x=1$, η $C_f$ έχει σημεία καμπής τα
\[
A(-1,\ln2)
\quad\text{και}\quad
M(1,\ln2).
\]
\end{erwthma}
\end{thema}
%=========== ΤΕΛΟΣ ΛΥΣΗΣ - 64ο ΘΕΜΑ Β ===========
